% Hafta 11 — Dış kodlama: en güçlü koruma, en büyük sınır
% Derleme: py -3.12 tools/latex/derle.py docs/week-11/assets/h11-04-dis-kodlama.tex
\documentclass[tikz,border=8pt]{standalone}
\usepackage{cen429-cizim}
\begin{document}
\begin{tikzpicture}[node distance=6mm and 22mm]
  \node[baslik] (b) at (0,0) {Dış kodlama: en güçlü koruma, en büyük sınır};
  \node[altbaslik] at (b.south west) {çıktı artık STANDART AES değildir};

  \node[kutu=30mm, below=16mm of b.south west, anchor=north west] (n1) {\kb{Girdi}};
  \node[uyarikutu=30mm, right=of n1] (n2) {\kb{$F^{-1}$}\\(dış kodlama)\\uygulamaya bağlar};
  \node[koyukutu=30mm, right=of n2] (n3) {\kb{Kodlanmış\\tablo ağı}\\AES turları};
  \node[uyarikutu=30mm, right=of n3] (n4) {\kb{G}\\(dış kodlama)};
  \node[kotukutu=34mm, right=of n4] (n5) {\kb{Çıktı}\\$G \circ \mathrm{AES} \circ F^{-1}$};
  \draw[ok, draw=soluk] (n1.east) -- (n2.west);
  \draw[ok, draw=soluk] (n2.east) -- (n3.west);
  \draw[ok, draw=soluk] (n3.east) -- (n4.west);
  \draw[ok, draw=soluk] (n4.east) -- (n5.west);

  \node[kotudip, text width=160mm, below=8mm of n5.south -| n3, anchor=north] {%
    Karşı taraf F ve G'yi bilmeden sonuç anlamsız $\rightarrow$ yalnız iki ucu da sizin olduğu protokolde.};
\end{tikzpicture}
\end{document}
